\documentclass[10pt,a4paper]{amsart}
\usepackage[margin=19mm]{geometry}
\usepackage{amsmath,amssymb,mathtools}
\usepackage[scr=rsfs,cal=euler]{mathalfa}
\usepackage{xfrac}
\usepackage[unicode,hidelinks,bookmarks=false]{hyperref}
\newcommand{\ie}{i.e.\ }
\newcommand{\cf}{cf.\ }
\newcommand{\resp}{respectively\ }
\newcommand{\Subsection}{\subsection}
\providecommand{\nobreakdash}{\nobreak\hbox{-}}
\setlength{\emergencystretch}{2em}
\multlinegap0pt
\usepackage{bm}
\usepackage{bbm}
\usepackage{dsfont}
\usepackage{xspace}
\usepackage{cancel}
\usepackage{mathtools}
\usepackage{amsthm}
\makeatletter
\@ifundefined{theorem}{\newtheorem{theorem}{Theorem}}{}
\@ifundefined{lemma}{\newtheorem{lemma}[theorem]{Lemma}}{}
\makeatother
\DeclareMathOperator*{\tr}{tr}
\DeclareMathOperator{\E}{\mathbb{E}}
\DeclareMathOperator{\Pb}{\mathbb{P}}
\DeclareMathOperator{\R}{\mathbb{R}}
\DeclarePairedDelimiter{\norm}{\lVert}{\rVert}
\DeclarePairedDelimiter{\paren}{(}{)}
\newcommand{\norml}[1]{\left\lVert#1\right\rVert}
\title[Optimal Subspace Embeddings: Resolving Nelson-Nguyen Conject]{Optimal Subspace Embeddings: Resolving Nelson-Nguyen Conjecture Up to Sub-Polylogarithmic Factors}
\author{Shabarish Chenakkod, Micha{\l} Derezi\'nski, Xiaoyu Dong}
\date{Source: arXiv:2508.14234v2 (CC BY 4.0)}
\begin{document}
\maketitle

\noindent\emph{Source and licence.} Optimal Subspace Embeddings: Resolving Nelson-Nguyen Conjecture Up to Sub-Polylogarithmic Factors. Version arXiv:2508.14234v2 (CC BY 4.0), distributed under the Creative Commons Attribution 4.0 International licence, as recorded in the arXiv licence metadata for this exact version and shown on the abstract page https://arxiv.org/abs/2508.14234v2. Full rights notice: licenses/arxiv-2508.14234-CC-BY-4.0.txt.\par\smallskip

\noindent\textit{Editorial scope.} Counted unit: the Lemma whose source label is \texttt{lem:iterdec} in the article (source0.tex lines 1498--1525, source label \texttt{lem:iterdec}), delivered together with the article's own complete original proof of it (source0.tex lines 1529--1770, one \texttt{proof} environment of 242 lines). No mathematics is written, repaired or re-derived: statement and proof are the article's own text line for line. The proof is detached from the statement in the source: the article states the result at lines 1498--1525 and prints its proof at lines 1529--1770, and the proof environment's own header names the statement, so the association is the article's own. Required context retained uncounted: lem:initialuniv (lines 1253--1271); lem:decoupfirstlayer (lines 1468--1476); lem:decoupsecondlayer (lines 1479--1487); lem:x1x2unorm (lines 2760--2765). Every definition, display and label of those lines is retained unchanged. Omitted from this package and not redistributed: 1--1252, 1272--1467, 1477--1478, 1488--1497, 1526--1528, 1771--2759, 2766--3148. Nothing inside a retained excerpt is omitted or altered: every retained body is delivered exactly as the article's lines are written, so no omission receipt of any kind accompanies this package. Citations: the retained excerpts carry no citation command at all, so no bibliography entry of the article is transcribed and no reference is redistributed. Reference adaptation: none was needed and none was made; every reference inside the delivered unit resolves inside it, so no unresolved reference or citation is produced. Printed slips kept literal: no printed word, symbol, hypothesis or number of the counted statement or of its proof was altered. Layout and class: the article's own class is \texttt{amsart} with packages including geometry, amsfonts, amssymb, amsmath, enumerate, mathtools, amsthm, hyperref, amsaddr, bbm, xcolor, biblatex, wasysym, mathrsfs; the wrapper is the lane's pinned \texttt{amsart} editorial wrapper at 10pt on A4, which restates the theorem-like environments the retained text needs and adds the article's own macro definitions that the retained text needs (\texttt{\textbackslash E}, \texttt{\textbackslash Pb}, \texttt{\textbackslash R}, \texttt{\textbackslash norm}, \texttt{\textbackslash norml}, \texttt{\textbackslash paren}, \texttt{\textbackslash tr}), with extra wrapper packages (bm, bbm, dsfont, xspace, cancel, mathtools). The delivered numbering is the unit's own. Assets: no retained span contains a figure environment, a graphics-inclusion command, a PDF-page inclusion, a table or a TikZ picture; no image file is copied into this package, the package declares no asset dependency and no image file is redistributed with it. Licence: version arXiv:2508.14234v2 of this article is distributed under the Creative Commons Attribution 4.0 International licence, as recorded in the arXiv licence metadata for this exact version and shown on the abstract page https://arxiv.org/abs/2508.14234v2. A full rights notice is delivered at licenses/arxiv-2508.14234-CC-BY-4.0.txt. Characters: every retained excerpt is pure ASCII, so the delivered engine, XeLaTeX with its default Latin Modern text font, reports no missing character.
\begin{lemma}[Initial Moment Estimate]\label{lem:initialuniv}
Let $S$ be an $m \times n$ matrix distributed according to the fully independent unscaled OSNAP distribution with parameter $p$ and standard independent summands $\{Z_{(l,\gamma)}\}_{(l,\gamma) \in \Xi}$. Let $\{(S_{\lambda},\{Z_{\lambda,(l,\gamma)}\}_{(l,\gamma) \in \Xi})\}_{\lambda=0,1,2,...}$ be independent copies of $(S,\{Z_{(l,\gamma)}\}_{(l,\gamma) \in \Xi})$. Let $U$ be an arbitrary $n \times d$ deterministic matrix such that $U^TU=I_d$.
Let $X=SU$, and $X_1,X_2,...$ be i.i.d. copies of $X$. Let $Y=X_{1}^TX_{2}$. For any $d \times d$ deterministic matrix $V$ and positive integer $r$, define
\begin{align*}
    R_{1,2r}(V)=(\sum\limits_{(l,\gamma) \in \Xi}(\norm{VX_1^TZ_{2,(l,\gamma)}U}_{2r}^{2r}))^{1/(2r)}
\end{align*}
and
\begin{align*}
    R_{2,2r}(V)=(\sqrt{pd+2r}) (\sum\limits_{(l,\gamma) \in \Xi}\frac{1}{d}\norm{Vu_l}_{l_2([d])}^{2r})^{1/(2r)}
\end{align*}
Then there exists a constant $c_{\ref{lem:initialuniv}}>1$ such that for any $n \times d$ deterministic matrix $V$, and $\log(d) \le q \le r$, we have
\begin{equation}\label{eq:initmom}
    \begin{aligned}
    &\norm{VY}_{2q} \\\le&  c_{\ref{lem:initialuniv}}(\sqrt{p\max\{m,q\}p\max\{d,q\}}\norm{V} + q^{2}R_{2,2r}(V) + q^{2}R_{1,2r}(V))
\end{aligned}
\end{equation}


\end{lemma}

\begin{lemma}[First Layer Decoupling]\label{lem:decoupfirstlayer}
Let $S$, $p$, $\{Z_{(l,\gamma)}\}_{(l,\gamma) \in \Xi}$, $\{(S_{\lambda},\{Z_{\lambda,(l,\gamma)}\}_{(l,\gamma) \in \Xi})\}_{\lambda=0,1,2,...}$, $U$, $X$, $X_1,X_2,...$, and $Y_1,Y_2,...$ be as in Lemma \ref{lem:initialuniv}.
Then for any $d \times d$ deterministic matrix $V$, we have
\begin{align*}
    &\norm{Y_1^TV^TVY_1}_{2q} \\ \le&
    \norm{V}^2 c_{\ref{lem:decoupfirstlayer}}(pmpd+q^2+(pm)^{1/(q)}q(pd+q))\\&+4\norm{X_3^TX_1V^TVX_1^TX_4}_{2q}
\end{align*}

\end{lemma}

\begin{lemma}[Second Layer Decoupling]\label{lem:decoupsecondlayer}
Let $S$, $p$, $\{Z_{(l,\gamma)}\}_{(l,\gamma) \in \Xi}$, $\{(S_{\lambda},\{Z_{\lambda,(l,\gamma)}\}_{(l,\gamma) \in \Xi})\}_{\lambda=0,1,2,...}$, $U$, $X$, $X_1,X_2,...$, and $Y_1,Y_2,...$ be as in Lemma \ref{lem:initialuniv}.
Then for any $d \times d$ deterministic matrix $V$, we have
\begin{align*}
    &\norm{X_3^TX_1V^TVX_1^TX_4}_{2q} \\ \le&
    \norm{V}^2 c_{\ref{lem:decoupsecondlayer}}(pmpd+pd \norm{Y_1}_{2q}+(pm)^{1/(q)}q(pd+q))\\&+4\norm{X_3^TX_5V^TVX_6^TX_4}_{2q}
\end{align*}

\end{lemma}

\begin{lemma}\label{lem:x1x2unorm}
    Let $S$ be an $m \times n$ matrix distributed according to the fully independent unscaled OSNAP distribution with parameter $p$. Let $U$ be an arbitrary $n \times d$ deterministic matrix such that $U^TU=\operatorname{Id}$. Let $X=SU$. Let $X_1$ and $X_2$ be i.i.d. copies of $X$. Let $u \in \R^d$ be a vector with $\norm{u^T}_{l_2([d])}=1$. Let $q \ge 1$ be an integer. Then there exists a constant $c_{\ref{lem:x1x2unorm}}$ such that
\begin{align*}
    \norm{\norm{X_1^TX_2u}_2}_{L_{2q}(\Pb)} \le c_{\ref{lem:x1x2unorm}}\paren*{\sqrt{pmpd}+ (pmd)^\frac{1}{q} (\sqrt{pdq} + q)}
\end{align*}
\end{lemma}

\begin{lemma}[Iteratively Improved Moment Estimate]\label{lem:iterdec}
Let $S$ be an $m \times n$ matrix distributed according to the fully independent unscaled OSNAP distribution with parameter $p$ and standard independent summands $\{Z_{(l,\gamma)}\}_{(l,\gamma) \in \Xi}$. Let $\{(S_{\lambda},\{Z_{\lambda,(l,\gamma)}\}_{(l,\gamma) \in \Xi})\}_{\lambda=0,1,2,...}$ be independent copies of $(S,\{Z_{(l,\gamma)}\}_{(l,\gamma) \in \Xi})$. Let $U$ be an arbitrary $n \times d$ deterministic matrix such that $U^TU=I_d$.
Let $X=SU$. Let $Y_i=X_{2i-1}^TX_{2i}$.
For any $d \times d$ deterministic matrix $V$ and positive integer $r$, let $R_{1,2r}(V)$ and $R_{2,2r}(V)$ be defined as in Lemma \ref{lem:initialuniv}.
Let
\begin{align*}
    K_{q,r}=&K(m,d,p,q,r)\\=&\left(p\max\{m,q\}p\max\{d,q\}+(pm)^{1/(q)}r(pd+r) \right)^{1/2}
\end{align*}
Then there exist constants $c_{\ref{lem:iterdec}.1}>1$ and $c_{\ref{lem:iterdec}.2}>1$ such that for any $q \ge \log(d)$, $\alpha_1>0,...,\alpha_7>0$, and $\kappa > c_{\ref{lem:iterdec}.1}$, the following statement is true.

Assume that
    $\alpha_3 \le 1$ and $\alpha_1+\alpha_3=1$
and for any $n \times d$ deterministic matrix $V$ and $r \ge q$, we have
\begin{equation}\label{eq:iterassum}
    \begin{aligned}
    &\norm{VY}_{2q} \\\le& \kappa \Bigg((K_{q,r}+\sqrt{pd\norm{Y}_{2q}})\norm{V} \\&+ \left(R_{1,2r}(V)+R_{2,2r}(V)\right)^{\alpha_1}K_{q,r}^{\alpha_2}\norm{V}^{\alpha_3} \norm{Y}_{2q}^{\alpha_4}q^{\alpha_5}\Bigg)
\end{aligned}
\end{equation}
then for any $n \times d$ deterministic matrix $V$ and $r \ge 2q$, we have
\begin{equation}\label{eq:iterconclu}
    \begin{aligned}
    &\norm{VY}_{4q} \\\le& c_{\ref{lem:iterdec}.2}\kappa \Bigg((K_{q,r}+\sqrt{pd\norm{Y}_{4q}})\norm{V} \\&+ \left(R_{1,2r}(V)+R_{2,2r}(V)\right)^{\alpha_1/2}K_{q,r}^{(\alpha_2+\min\{\frac{1}{2},\alpha_1\})/2}\norm{V}^{\alpha_3+\alpha_1/2} \norm{Y}_{4q}^{(\alpha_4+\max\{\frac{1}{2},\alpha_3\})/2}q^{\alpha_5/2}\Bigg)
\end{aligned}
\end{equation}



\end{lemma}

\begin{proof}[Proof of Lemma \ref{lem:iterdec}]


Assume that (\ref{eq:iterassum}) holds for some $\alpha_1>0,...,\alpha_5>0$.

\textbf{Step 1: decoupling.}


Consider an arbitrary $n \times d$ deterministic matrix $V$ and $r \ge q$. If $\norm{Y}_{4q} \le K_{q,r}$, then we have
\begin{align*}
    &\norm{VY}_{4q} \\\le& 
    \norm{V}\norm{Y}_{4q}
    \\ \le & K_{q,2r}\norm{V}
    \\ \le &\kappa \Bigg((K_{q,r}+\sqrt{pd\norm{Y}_{4q}})\norm{V} \\&+ \left(R_{1,2r}(V)+R_{2,2r}(V)\right)^{\alpha_1/2}K_{q,r}^{(\alpha_2+\min\{\frac{1}{2},\alpha_1\})/2}\norm{V}^{\alpha_3+\alpha_1/2} \norm{Y}_{4q}^{(\alpha_4+\max\{\frac{1}{2},\alpha_3\})/2}q^{\alpha_5/2}\Bigg)
\end{align*}
where we use $\kappa>1$, and therefore the desired statement (\ref{eq:iterconclu}) holds automatically.

Therefore, it suffices to only consider the case when $\norm{Y}_{4q} > K_{q,r}$.


We have
\begin{align*}
    \norm{VY}_{4q}=&(\E\tr((Y^TV^TVY)^{2q}))^{\frac{1}{4q}}
\\=&\norm{Y^TV^TVY}_{2q}^{1/2}
\end{align*}

By Lemma \ref{lem:decoupfirstlayer} and Lemma \ref{lem:decoupsecondlayer}, we know that
\begin{align*}
    \norm{Y_1^TV^TVY_1}_{2q} \le& c_1 K^2 \norm{V}^2 + 4pd\norm{V}^2 \norm{Y_1}_{2q} + 16 \norm{X_3^TX_5V^TVX_6^TX_4}_{2q}
    \\=& c_1 K^4 \norm{V}^2 + 4pd\norm{V}^2 \norm{Y_1}_{2q}+ 16 \norm{Y_3^TV^TVY_4}_{2q}
    \\ \le & c_1 K^4 \norm{V}^2 + 4pd\norm{V}^2 \norm{Y_1}_{4q}+ 16 \norm{Y_3^TV^TVY_4}_{2q}
\end{align*}



\textbf{Step 2: conditioning on $W=Y_3^TV^TV$.}

Now we will condition on $W=Y_3^TV^TV$ and use our assumption to bound $\norm{WY_4}_{2q}$ for fixed $W$.

Since $r \ge 2q \ge q$, by our assumption (\ref{eq:iterassum}), we have
    \begin{align*}
    \norm{WY}_{2q} \le & \kappa \Bigg((K_{q,r}+\sqrt{pd\norm{Y}_{2q}})\norm{W} \\&+ \left(R_{1,2r}(W)+R_{2,2r}(W)\right)^{\alpha_1}K_{q,r}^{\alpha_2}\norm{W}^{\alpha_3} \norm{Y}_{2q} ^{\alpha_4}q^{\alpha_5}\Bigg)
    \\\le & \kappa \Bigg((K_{q,r}+\sqrt{pd\norm{Y}_{4q}})\norm{W} \\&+ \left(R_{1,2r}(W)+R_{2,2r}(W)\right)^{\alpha_1}K_{q,r}^{\alpha_2}\norm{W}^{\alpha_3} \norm{Y}_{4q} ^{\alpha_4}q^{\alpha_5}\Bigg)
\end{align*}
where we use the fact $\norm{Y}_{2q} \le \norm{Y}_{4q}$ by H\"older inequality in $L_q(S_q^d)$.

Since $q \ge \log(d)$, we have $\norm{\norm{W}}_{L_{2q}(\Pb)}  \le d^{\frac{1}{\log d}} \norm{W}_{2q}= e \norm{W}_{2q}$. Therefore, by triangle inequality for $\norm{\cdot}_{L_{4q}(\Pb)}$, we have
    \begin{align*}
    &\norm{Y_3^TV^TVY}_{2q} \\\le& e\kappa \Bigg((K_{q,r}+\sqrt{pd\norm{Y}_{4q}})\norm{Y_3^TV^TV}_{2q} \\&+ \norml{\left(R_{1,2r}(Y_3^TV^TV)+R_{2,2r}(Y_3^TV^TV)\right)^{\alpha_1}K_{q,r}^{\alpha_2}\norm{Y_3^TV^TV}^{\alpha_3} \norm{Y}_{4q} ^{\alpha_4}q^{\alpha_5}}_{L_{2q}(\Pb)}\Bigg)
\end{align*}


We first estimate the first term
\begin{align*}
    (K_{q,r}+\sqrt{pd\norm{Y}_{4q}})\norm{Y_3^TV^TV}_{2q}
\end{align*}

By using assumption (\ref{eq:iterassum}) again, we have
\begin{align*}
    \norm{Y_3^TV^TV}_{2q}=&\norm{V^TVY_3}_{2q} \\\le& \kappa \Bigg((K_{q,r}+\sqrt{pd\norm{Y}_{2q}})\norm{V^TV} \\&+ \left(R_{1,2r}(V^TV)+R_{2,2r}(V^TV)\right)^{\alpha_1}K_{q,r}^{\alpha_2}\norm{V^TV}^{\alpha_3} \norm{Y}_{2q}^{\alpha_4}q^{\alpha_5}\Bigg)
    \\\le& \kappa \Bigg((K_{q,r}+\sqrt{pd\norm{Y}_{4q}})\norm{V}^2 \\&+ \left(R_{1,2r}(V)+R_{2,2r}(V)\right)^{\alpha_1}K_{q,r}^{\alpha_2}\norm{V}^{2\alpha_3+\alpha_1} \norm{Y}_{4q}^{\alpha_4}q^{\alpha_5}\Bigg)
\end{align*}
where we use the fact that
\begin{align*}
    R_{1,2r}(V^TV)=&(\sum\limits_{(l,\gamma) \in \Xi}(\norm{V^TVX_1^TZ_{2,(l,\gamma)}U}_{2r}^{2r}))^{1/(2r)}
    \\ \le &(\sum\limits_{(l,\gamma) \in \Xi}((\norm{V}\norm{VX_1^TZ_{2,(l,\gamma)}U}_{2r})^{2r}))^{1/(2r)}
    \\ = &\norm{V}(\sum\limits_{(l,\gamma) \in \Xi}(\norm{VX_1^TZ_{2,(l,\gamma)}U}_{2r}^{2r}))^{1/(2r)}
    \\=& \norm{V}R_{1,2r}(V)
\end{align*}
and
\begin{align*}
    R_{2,2r}(V^TV) \le & \norm{V}R_{2,2r}(V)
\end{align*}
by the same calculation.

Next, we bound the second term
\begin{align*}
    \norml{\left(R_{1,2r}(Y_3^TV^TV)+R_{2,2r}(Y_3^TV^TV)\right)^{\alpha_1}K_{q,r}^{\alpha_2}\norm{Y_3^TV^TV}^{\alpha_3} \norm{Y}_{4q} ^{\alpha_4}q^{\alpha_5}}_{L_{2q}(\Pb)}
\end{align*}

We have
\begin{align*}
    &\norml{\left(R_{1,2r}(Y_3^TV^TV)+R_{2,2r}(Y_3^TV^TV)\right)^{\alpha_1}K_{q,r}^{\alpha_2}\norm{Y_3^TV^TV}^{\alpha_3} \norm{Y}_{4q} ^{\alpha_4}q^{\alpha_5}}_{L_{2q}(\Pb)}
    \\=&K_{q,r}^{\alpha_2}\norm{Y}_{4q}^{\alpha_4}q^{\alpha_5}\norml{\left(R_{1,2r}(Y_3^TV^TV)+R_{2,2r}(Y_3^TV^TV)\right)^{\alpha_1}\norm{Y_3^TV^TV}^{\alpha_3}  }_{L_{2q}(\Pb)}
    \\ \le &K_{q,r}^{\alpha_2}\norm{Y}_{4q}^{\alpha_4}q^{\alpha_5}\norml{\left(R_{1,2r}(Y_3^TV^TV)+R_{2,2r}(Y_3^TV^TV)\right)^{\alpha_1}\norm{Y_3}^{\alpha_3} \norm{V}^{2 \alpha_3}  }_{L_{2q}(\Pb)}
    \\ \le &K_{q,r}^{\alpha_2}\norm{Y}_{4q}^{\alpha_4}\norm{V}^{2 \alpha_3}q^{\alpha_5}\norml{\left(R_{1,2r}(Y_3^TV^TV)+R_{2,2r}(Y_3^TV^TV)\right)^{\alpha_1}\norm{Y_3}^{\alpha_3}   }_{L_{2q}(\Pb)}
    \\ \le &K_{q,r}^{\alpha_2}\norm{Y}_{4q}^{\alpha_4}\norm{V}^{2 \alpha_3}q^{\alpha_5}\norml{\left(R_{1,2r}(Y_3^TV^TV)+R_{2,2r}(Y_3^TV^TV)\right)^{\alpha_1}}_{L_{4q}(\Pb)}\norml{\norm{Y_3}^{\alpha_3}   }_{L_{4q}(\Pb)}
\end{align*}
where in the last step, we use Holder's inequality.

We also have two observations. Since $r \ge 2q$ and $\alpha_1 \le 1$, we have
\begin{align*}
    &\norml{\left(R_{1,2r}(Y_3^TV^TV)+R_{2,2r}(Y_3^TV^TV)\right)^{\alpha_1}}_{L_{4q}(\Pb)}
    \\ \le & \norml{\left(R_{1,2r}(Y_3^TV^TV)+R_{2,2r}(Y_3^TV^TV)\right)^{\alpha_1}}_{L_{2r}(\Pb)}
    \\ \le & \norml{\left(R_{1,2r}(Y_3^TV^TV)+R_{2,2r}(Y_3^TV^TV)\right)}_{L_{2r}(\Pb)}^{\alpha_1}
\end{align*}
And since $\alpha_3 \le 1$, we have
\begin{align*}
    \norml{\norm{Y_3}^{\alpha_3}   }_{L_{4q}(\Pb)} \le c_2\norm{Y_3}_{4q}^{\alpha_3}
\end{align*}
for some constant $c_2$ by the fact that $q \ge \log(d)$ and Jensen's inequality.

To estimate $\norml{\left(R_{1,2r}(Y_3^TV^TV)+R_{2,2r}(Y_3^TV^TV)\right)}_{L_{2r}(\Pb)}^{\alpha_1}$, we need the following lemma. This step is main reason why this iterative method can improve the estimate on $\norm{Y}_{4q}$. Intuitively, in this step, we can take out a copy of $Y$ from $R_{1,2r}(Y_3^TV^TV)+R_{2,2r}(Y_3^TV^TV)$ and replace it by the correct factor $K_{q,r}$.

\begin{lemma}\label{lem:Rsimplify}
Let $R_{1,2r}(V), R_{2,2r}(V)$ and $K_{q,r}$ be as in Lemma \ref{lem:iterdec} with $r \ge2q$, then we have
\begin{align*}
    \norm{R_{1,2r}(Y_3^TV^TV)}_{L_{2r}(\Pb)}  
    \le &c_{\ref{lem:Rsimplify}}(R_{1,2r}(V))\norm{V}K_{q,r}
\end{align*}
and
\begin{align*}
    \norm{R_{2,2r}(Y_3^TV^TV)}_{L_{2r}(\Pb)}  
    \le &c_{\ref{lem:Rsimplify}}(R_{2,2r}(V))\norm{V}K_{q,r}
\end{align*}
\end{lemma}

\begin{proof}
By definition, we have
\begin{align*}
    R_{1,2r}(Y_3^TV^TV)=(\sum\limits_{(l,\gamma) \in \Xi}(\E_{X_1,X_2} \tr(|Y_3^TV^TVX_1^TZ_{2,(l,\gamma)}U|^{2r})))^{1/(2r)}
\end{align*}

Therefore, we have
\begin{align*}
    \norm{R_{1,2r}(Y_3^TV^TV)}_{L_{2r}(\Pb)}^{2r}  
    = & \E_{Y_3} \sum\limits_{(l,\gamma) \in \Xi}(\E_{X_1,X_2} \tr(|Y_3^TV^TVX_1^TZ_{2,(l,\gamma)}U|^{2r}))
    \\= &  \sum\limits_{(l,\gamma) \in \Xi}(\E_{X_1,X_2} \E_{Y_3}\tr(|Y_3^TV^TVX_1^TZ_{2,(l,\gamma)}U|^{2r}))
\end{align*}





Recall that $Z_{k,(l,\gamma)}=\xi_{k,(l,\gamma)}e_{\mu_{k,(l,\gamma)}}e_l^T$. Also, define the vector $u_l$ such that $u_l^T$ is the $l$th row of the matrix $U$. Therefore, we observe that
\begin{align*}
    &\tr(|Y_3^TV^TVX_1^TZ_{2,\eta_2}U|^{2r})
\\=&\tr(|Y_3^TV^TVX_1^T\xi_{2,\eta_2}e_{\mu_{2,\eta_2}}e_{\eta_2(1)}^TU|^{2r})
\\=&\tr(|Y_3^TV^TVX_1^T\xi_{2,\eta_2}e_{\mu_{2,\eta_2}}u_{\eta_2(1)}^T|^{2r})
\\=&\frac 1 d\norm{Y_3^TV^TVX_1^Te_{\mu_{2,\eta_2}}}_{l_2([d])}^{2r}\norm{u_{\eta_2(1)}}_{l_2([d])}^{2r}
\end{align*}

Therefore, we have
\begin{align*}
    &E_{Y_3}\tr(|Y_3^TV^TVX_1^TZ_{2,\eta_2}U|^{2r})
\\=&\frac 1 dE_{Y_3}\norm{Y_3^TV^TVX_1^Te_{\mu_{2,\eta_2}}}_{l_2([d])}^{2r}\norm{u_{\eta_2(1)}}_{l_2([d])}^{2r}
\\=&\frac 1 d\norm{u_{\eta_2(1)}}_{l_2([d])}^{2r}E_{Y_3}\norm{Y_3^TV^TVX_1^Te_{\mu_{2,\eta_2}}}_{l_2([d])}^{2r}\\=&\frac 1 d\norm{u_{\eta_2(1)}}_{l_2([d])}^{2r}E_{X_6}\E_{X_5}\norm{X_5^TX_6V^TVX_1^Te_{\mu_{2,\eta_2}}}_{l_2([d])}^{2r}
\end{align*}

By Lemma \ref{lem:x1x2unorm}, we have
\begin{align*}
&\E_{X_5,X_6}\norm{X_5^TX_6V^TVX_1^Te_{\mu_{2,\eta_2}}}_{l_2([d])}^{2r} \\\le& \left(\norm{V^TVX_1^Te_{\mu_{2,\eta_2}}}_{l_2([m])}c_{\ref{lem:x1x2unorm}}\paren*{\sqrt{pmpd}+ (pmd)^\frac{1}{q} (\sqrt{pd\cdot 2r} + 2r)}\right)^{2r}
\\\le&\norm{V^TVX_1^Te_{\mu_{2,\eta_2}}}_{l_2([m])}^{4q}c_1K_{q,r}^{2r}
\end{align*}
where we use
\begin{align*}
    \sqrt{pmpd}+ (pmd)^\frac{1}{q} (\sqrt{pdq} + q) \le c_1K_{q,r}
\end{align*}
for some constant $c_1$ just by definition of $K_{q,r}$.

Therefore, we have
\begin{align*}
    &E_{Y_3}\tr(|Y_3^TV^TVX_1^TZ_{2,\eta_2}U|^{2r})
\\=&\frac 1 d\norm{u_{\eta_1(1)}}_{l_2([d])}^{2r}\norm{V^TVX_1^Te_{\mu_{2,\eta_2}}}_{l_2([m])}^{2r}(c_{1}K_{q,r})^{2r}
\\ \le & \frac 1 d\norm{V}^{2r} \norm{u_{\eta_2(1)}}_{l_2([d])}^{2r}\norm{VX_1^Te_{\mu_{2,\eta_2}}}_{l_2([m])}^{2r}(c_{1}K_{q,r})^{2r}
\\=&\norm{V}^{2r} (c_{1}K_{q,r})^{2r} \tr(|VX_1^Te_{\mu_{2,\eta_2}}u_{\eta_2(1)}|^{2r})
\end{align*}


Therefore, we have
\begin{align*}
    \norm{R_{1,2r}(Y_3^TV^TV)}_{L_{2r}(\Pb)}^{2r}  
    = &  \sum\limits_{(l,\gamma) \in \Xi}(\E_{X_1,X_2} \E_{Y_3}\tr(|Y_3^TV^TVX_1^TZ_{2,(l,\gamma)}U|^{2r}))
    \\ \le &\sum\limits_{(l,\gamma) \in \Xi}( \norm{V}^{2r} (c_{1}K_{q,r})^{2r} \E_{X_1,X_2}\tr(|VX_1^Te_{\mu_{2,\eta_2}}u_{\eta_2(1)}|^{2r}))
\end{align*}

Recall that, by definition, we have
\begin{align*}
    (R_{1,2q}(V))^{2r}=&\sum\limits_{(l,\gamma) \in \Xi}(\norm{VX_1^TZ_{2,(l,\gamma)}U}_{2r}^{2r})
    \\=&\sum\limits_{(l,\gamma) \in \Xi}\E_{X_1,X_2}\tr(|VX_1^Te_{\mu_{2,\eta_2}}u_{\eta_2(1)}|^{2r})
\end{align*}

Taking $2r$-th root of both sides, we have
\begin{align*}
    &\norm{R_{1,2r}(Y_3^TV^TV)}_{L_{2r}(\Pb)}  
    \\ \le &\norm{V} (c_{1}K_{q,r})R_{1,2r}(V)
\end{align*}
which is what we want.

The proof for the inequality for $\norm{R_{2,2r}(Y_3^TV^TV)}_{L_{2r}(\Pb)}  $ follows exactly the same argument.



\end{proof}

Therefore, by Lemma \ref{lem:Rsimplify} and triangle inequality, we have
\begin{align*}
    &\norml{\left(R_{1,2r}(Y_3^TV^TV)+R_{2,2r}(Y_3^TV^TV)\right)^{\alpha_1}K_{q,r}^{\alpha_2}\norm{Y_3^TV^TV}^{\alpha_3} \norm{Y}_{4q} ^{\alpha_4}q^{\alpha_5}}_{L_{2q}(\Pb)}
    \\ \le &c_2K_{q,r}^{\alpha_2}\norm{Y}_{4q}^{\alpha_4+\alpha_3}\norm{V}^{2 \alpha_3}q^{\alpha_5}\norml{\left(R_{1,2r}(Y_3^TV^TV)+R_{2,2r}(Y_3^TV^TV)\right)}_{L_{2r}(\Pb)}^{\alpha_1}
    \\ \le & c_2 K_{q,r}^{\alpha_2+\alpha_1}\norm{Y}_{4q}^{\alpha_4+\alpha_3}\norm{V}^{2 \alpha_3+\alpha_1}q^{\alpha_5}(\left(R_{1,2r}(V)+R_{2,2r}(V)\right)^{\alpha_1})
\end{align*}

Therefore, we have
    \begin{align*}
    &\norm{Y_3^TV^TVY}_{2q} \\\le& e\kappa \Bigg((K_{q,r}+\sqrt{pd\norm{Y}_{4q}})\norm{Y_3^TV^TV}_{2q} \\&+ \norml{\left(R_{1,2r}(Y_3^TV^TV)+R_{2,2r}(Y_3^TV^TV)\right)^{\alpha_1}K_{q,r}^{\alpha_2}\norm{Y_3^TV^TV}^{\alpha_3} \norm{Y}_{4q} ^{\alpha_4}q^{\alpha_5}}_{L_{2q}(\Pb)}\Bigg)
    \\ \le & e\kappa \Bigg((K_{q,r}+\sqrt{pd\norm{Y}_{4q}}) \\& \cdot \kappa \Bigg((K_{q,r}+\sqrt{pd\norm{Y}_{4q}})\norm{V}^2 + \left(R_{1,2r}(V)+R_{2,2r}(V)\right)^{\alpha_1}K_{q,r}^{\alpha_2}\norm{V}^{2\alpha_3+\alpha_1} \norm{Y}_{4q}^{\alpha_4}q^{\alpha_5}\Bigg) \\&+ c_2 K_{q,r}^{\alpha_2+\alpha_1}\norm{Y}_{4q}^{\alpha_4+\alpha_3}\norm{V}^{2 \alpha_3+\alpha_1}q^{\alpha_5}(\left(R_{1,2r}(V)+R_{2,2r}(V)\right)^{\alpha_1})\Bigg)
    \\ = & e\kappa^2 (K_{q,r}+\sqrt{pd\norm{Y}_{4q}})^2 \norm{V}^2 \\&+ \kappa^2 (K_{q,r}+\sqrt{pd\norm{Y}_{4q}})\left(R_{1,2r}(V)+R_{2,2r}(V)\right)^{\alpha_1}K_{q,r}^{\alpha_2}\norm{V}^{2\alpha_3+\alpha_1} \norm{Y}_{4q}^{\alpha_4}q^{\alpha_5} \\& + c_3 \kappa K_{q,r}^{\alpha_2+\alpha_1}\norm{Y}_{4q}^{\alpha_4+\alpha_3}\norm{V}^{2 \alpha_3+\alpha_1}q^{\alpha_5}(\left(R_{1,2r}(V)+R_{2,2r}(V)\right)^{\alpha_1})
\\ \le  & e\kappa^2 (K_{q,r}+\sqrt{pd\norm{Y}_{4q}})^2 \norm{V}^2 \\&+ \kappa^2 \left(R_{1,2r}(V)+R_{2,2r}(V)\right)^{\alpha_1}K_{q,r}^{\alpha_2+\frac{1}{2}}\norm{V}^{2\alpha_3+\alpha_1} \norm{Y}_{4q}^{\alpha_4+\frac{1}{2}}q^{\alpha_5} \\& + c_3 \kappa K_{q,r}^{\alpha_2+\alpha_1}\norm{Y}_{4q}^{\alpha_4+\alpha_3}\norm{V}^{2 \alpha_3+\alpha_1}q^{\alpha_5}(\left(R_{1,2r}(V)+R_{2,2r}(V)\right)^{\alpha_1})
\\ \le  & c_3\kappa^2 \big((K_{q,r}+\sqrt{pd\norm{Y}_{4q}})^2 \norm{V}^2 \\&+ \left(R_{1,2r}(V)+R_{2,2r}(V)\right)^{\alpha_1}K_{q,r}^{\alpha_2+\min\{\frac{1}{2},\alpha_1\}}\norm{V}^{2\alpha_3+\alpha_1} \norm{Y}_{4q}^{\alpha_4+\max\{\frac{1}{2},\alpha_3\}}q^{\alpha_5} \big)
\end{align*}
for some constant $c_3>0$,
where we use the assumption that $\norm{Y}_{4q}>K_{q,r}$ and the fact that $pd \le K_{q,r}$ to get
\begin{align*}
    K_{q,r}+\sqrt{pd\norm{Y}_{4q}} \le (\sqrt{K_{q,r}} + \sqrt{pd})\sqrt{\norm{Y}_{4q}} \le 2 \sqrt{K_{q,r}}\sqrt{\norm{Y}_{4q}}
\end{align*}
and we combine this with the fact that $\alpha_1+\alpha_3=1$ to conclude that the factor $\max \{ K_{q,r}^{\frac{1}{2}} \norm{Y}_{4q}^{\frac{1}{2}}, K_{q,r}^{\alpha_1} \norm{Y}_{4q}^{\alpha_3} \}$
is the first term when $\alpha_3<1/2 \text{ and } \alpha_1 >1/2$ and the second term otherwise. 


\textbf{Step 3: final estimates.}

Therefore, we have
\begin{align*}
    &\norm{Y_1^TV^TVY_1}_{2q}  \\\le & c_1 K_{q,r}^4 \norm{V}^2 + 4pd\norm{V}^2 \norm{Y_1}_{4q}+ 16 \norm{Y_3^TV^TVY_4}_{2q}
    \\ \le & c_4 \kappa^2 \big((K_{q,r}+\sqrt{pd\norm{Y}_{4q}})^2 \norm{V}^2 \\&+ \left(R_{1,2r}(V)+R_{2,2r}(V)\right)^{\alpha_1}K_{q,r}^{\alpha_2+\min\{\frac{1}{2},\alpha_1\}}\norm{V}^{2\alpha_3+\alpha_1} \norm{Y}_{4q}^{\alpha_4+\max\{\frac{1}{2},\alpha_3\}}q^{\alpha_5} \big)
\end{align*}
for some constant $c_4$.

Therefore, we have
\begin{align*}
    &\norm{VY}_{4q}
\\=&\norm{Y^TV^TVY}_{2q}^{1/2}
 \\ \le & c_5 \kappa \big((K_{q,r}+\sqrt{pd\norm{Y}_{4q}}) \norm{V} \\&+ \left(R_{1,2r}(V)+R_{2,2r}(V)\right)^{\alpha_1/2}K_{q,r}^{(\alpha_2+\min\{\frac{1}{2},\alpha_1\})/2}\norm{V}^{\alpha_3+\alpha_1/2} \norm{Y}_{4q}^{(\alpha_4+\max\{\frac{1}{2},\alpha_3\})/2}q^{\alpha_5/2} \big)
\end{align*}
for some constant $c_5$.






\end{proof}

\end{document}
